\documentclass{syllabus}

\begin{document}

\thispagestyle{fancy}
\date{Spring 2026}
\author{}
\title{PHY 4300: {Electricity and Magnetism}}
\maketitle

\meeting{Lecture: & MWF, 9:15-10:20 AM }
\officehours{TR: & 11:30 AM - 12:30 PM}
\textbook{Introduction to Electrodynamics (5/e)}{Griffiths}{Cambridge}{2024}
\prerequisites{PHY 2210: & (General Physics II) \\ & MTH 2020: & (Dif. eq.) \\ & MTH 2120: & (Vec. calc., co-enrolled)}

\section*{General Information}

\infotable
\hspace{\fill}
\includegraphics[height=2.25in]{Assets/electrodynamicsCover.png}
\hspace{.25in}





\section*{Course overview}

The study of the electric and magnetic effects of charges and currents leading to a presentation of Maxwell's equations and including such topics as electrostatic fields, electrostatic and magnetic energy, and potential theory.

Electricity and magnetism, in many ways, is THE field theory of physics. Many mathematical methods were developed especially for it: vector algebra, vector calculus, the theory of partial differential equations. E\&M, however, is not the only field theory. These methods have uses in many other fields, and a focus of this course will be in understanding those methods so that you can use them if you end up in one of those fields (elasticity theory, fluid dynamics, quantum field theory, etc.). 

PHY 4300 Electricity \& Magnetism is a study of the electric and magnetic effects of charges and currents, leading to a presentation of Maxwell’s equations and the understanding of light as an electromagnetic wave. A special approach, based on the fundamental physics of relativity and due to Purcell, is taken to understand fields in different frames of reference. Mathematical tools, notably vector calculus, are introduced in parallel with appropriate physics topics to emphasize their motivation and use. Topics include: electrostatic fields, electric potential, conductors, electric currents and circuits, fields of moving charges, magnetic fields, electromagnetic induction, Maxwell’s equations, and electromagnetic waves and radiation. 



\subsection*{Student learning outcomes} 

\noindent
Upon successful completion of this course, the student will be able to:
\begin{enumerate}[label=\textcolor{CarthageDark}{\bfseries\arabic*.}]
    \item Define and apply core principles of electromagnetism in the following areas: electrostatic fields, electric potential, conductors, electric currents and circuits, fields of moving charges, magnetic fields, electromagnetic induction, Maxwell’s equations, and electromagnetic waves/radiation.
    \item Solve canonical electromagnetism problems in upper-level undergraduate physics
\item Demonstrate understanding and insight into the fundamental principles of electromagnetism.
\item Apply basic numerical analysis to Poisson's equation and the visualization of scalar and vector fields.

\end{enumerate}


\section{Course components}
\paragraph*{Readings}
and their schedule are shown on the Course Schedule. Griffiths is a classic introductory text and is quite readable. Electricity and magnetism is a challenging subject and it is essential that you read the assigned text material before class and come prepared to discuss the assigned reading.

\paragraph{Homework} 
Homework exercises will be assigned on a daily basis and due once a week. You are encouraged to discuss them with fellow students and work in small groups: However, rote copying of someone else’s homework is prohibited. Late work will be accepted until solutions are posted ($\sim$1 week after due date) and will be penalized 30\%. 

\paragraph{Quizzes} 
Short quizzes on the content of each homework set will be administered on the day that homework set is due. Quizzes are mostly on a pass/fail system and are designed to test individual understanding of the material. 

\paragraph{Computational assignments} 
Over the course of the semester, you will create a numerical solver for Poisson's equation and for the curl-type Maxwell's equations. This code will be assessed as part of the homework, but at major intervals, the code and results thereof will be submitted as a project. 



\paragraph{Exams} 
Two cumulative mid-term exams will be administered during the semester. The format will be similar to homework problems (thought-provoking questions inviting you to explore your understanding) but suitably simplified to fit in a shorter time window. A comprehensive final exam will take place during the final exam period. You will be able to use a single handwritten notes sheet during the exams.

\clearpage

\section*{Grading Scheme}
\subsection*{Component Weights}

\noindent\begin{tabular}{lr}
\textbf{Component} & \textbf{Weight}       \\ \hline
\textbf{Participation}      & 0\%  \\
\textbf{Homework}           & 20\%       \\
\textbf{Quizzes}            & 10\%       \\
\textbf{Exams}              & 50\% \\
\textbf{Comp. projects}       & 20\%\\
\end{tabular}%



\lettergrades{a) number of extra credit assignments adequately completed, b) a pattern of seeking help in the course at office hours, c) number of missing homework assignments.}


\tipsforsuccess{upper}

\section*{Policies}
\subsection*{General course policies}
Expect a response within one business day for emails sent to the instructor. To ensure a prompt response, format the subject line as: [PHY 4300] *insert subject here*

All required material will be covered in lecture, but can also be found in the parallel readings, so in the case of absence students should consult the scheduled readings. Students found sleeping will be brought to the attention of the class for public ridicule and no effort made to wake them up.

Students are expected to arrive and leave on time. I will attempt to arrive five minutes early and be available for at least five minutes after class.

\latework{Easy}
\collegepolicies

\end{document}
